\documentclass{article}

\usepackage[utf8]{inputenc}
\usepackage{amsmath, esint}
\usepackage{wasysym}
\usepackage{qrcode}
\usepackage[colorlinks]{hyperref}
\usepackage{lmodern}
\usepackage{graphicx}
\usepackage{xcolor}
\usepackage[left=2cm, top=3cm, right=2cm]{geometry}
\usepackage{minted}
\usepackage{booktabs}
\usepackage{svg}
\usepackage{xcolor}
\usepackage{tabularx}
\definecolor{LightGray}{gray}{0.975}
\hypersetup{
  colorlinks,
  linkcolor=blue,
  urlcolor=blue
}

\title{ML101 \\ Lab 02: A `gentle' Introduction to Linear Model Selection \& Regularization.}
\author{Andrés Oswaldo Calderón Romero, Ph.D.}
\date{\today}

\begin{document}

\maketitle

\section{Introduction}
\label{sec:intro}

While Ordinary Least Squares (OLS) linear regression offers an intuitive baseline for predictive modeling, its performance degrades when confronting high-dimensional feature spaces, multicollinearity, or scarce observations relative to predictors ($p \gg n$). In such regimes, unconstrained least squares estimates suffer from high variance, causing severe overfitting and poor out-of-sample generalization. To mitigate these shortcomings, modern statistical learning relies on model selection and regularization paradigms: discrete subset selection (best subset, forward, and backward stepwise) isolates parsimonious predictor sets, shrinkage methods like Ridge ($\ell_2$) and Lasso ($\ell_1$) constrain coefficients to trade small amounts of bias for major variance reductions, and dimension reduction approaches like Principal Components Analysis project data onto orthogonal directions of maximal variance.

Building upon the foundations presented in Chapter~6 of \textit{An Introduction to Statistical Learning} (ISLRv2), this lab focuses on implementing, diagnosing, and contrasting these workflows on real-world empirical data. By operationalizing these techniques outside of the native R ecosystem (e.g., in Python), students will explore practical data pitfalls and evaluate the critical trade-offs among predictive accuracy, computational tractability, and model interpretability.

% \section{Additional Resources}
% \label{sec:resources}

\section{Lab: Linear Regression}
\label{sec:lab}

In this lab, we will work through the exercises in Chapter 6 of ISLRv2\footnote{James, G., Witten, D., Hastie, T., \& Tibshirani, R. (2023). \textit{An Introduction to Statistical Learning: with Applications in R} (2nd ed.). Springer. \url{https://www.statlearning.com/}}. On the section \textit{``6.5 Lab: Linear Models and Regularization Methods''} (page 267), we will explorer subsections:

\begin{itemize}
  \item 6.5.1 Subset Selection Methods,
  \item 6.5.2 Ridge Regression and the Lasso and
  \item 6.5.3 PCR and PLS Regression
\end{itemize}

You do not have to include any evidence in the lab report but you have to get familiar with the tools available to complete the autonomous work in next section.

\section{Autonomous Work}
\label{sec:work}

So, now we are able to apply selection and regularizatio methods to our data.  You can revisit the dataset you used in Lab 01 or pick a new one from the UCI Machine Learning Repository.  Indeed, we encorage you to explore your own data in this case.  Once you confirm what dataset you will use, it is expected that you apply what you learnt in section~\ref{sec:lab} in your dataset.  In detail, you will apply:

\begin{enumerate}
  \item Best Subset Selection.
  \item Forward Stepwise Selection.
  \item Backward Stepwise Selection.
  \item Ridge Regression.
  \item Lasso Regression.
  \item Principal components analysis.
\end{enumerate}

Do not worry if any of the techniques do not work or give you awful results.  The main idea is to known the available tools and how you manage any pitfall with the data.  You must use another programming language different to R.


\section{Deliverables}
Each group must submit a well-structured PDF report replicating the application of the methods proposed in Section~\ref{sec:work}.  It is expected that you discuss any relevant finding or particular remark. Submissions must be uploaded via the corresponding Brightspace\texttrademark{} assignment link within \textbf{21 days} of this lab session.

\vspace{5mm}
Happy Hacking! \includesvg[width=4mm]{figures/sunglasses}

\end{document}

